%MIT OpenCourseWare: https://ocw.mit.edu
%18.S190 Introduction to Metric Spaces, Independent Activities Period (IAP) 2023
%License: Creative Commons BY-NC-SA 
%For information about citing these materials or our Terms of Use, visit: https://ocw.mit.edu/terms.

\subsection*{Compact Metric Spaces}

Last time, we showed that a set in $\R^n$ is sequentially compact if and only if it is topologically compact, by showing 
\[
\text{sequentially compact} \iff \text{closed and bounded} \stackrel{\text{Heine-Borel}}{\iff} \text{topologically compact}.
\]
However, by the previous remark, we don't have Heine-Borel for arbitrary metric spaces. Which begs the question: is sequentially compact the same as topologically compact in metric spaces? The answer is yes. To prove this, we first show a handful of preliminary results.

\begin{lemma}[Lebesgue Number Lemma]
Let $(X,d)$ be a sequentially compact metric space and $\{U_i\}_{i\in I}$ be an open cover of $X$. Then, there exists an $r>0$ such that for each $x\in X$, $B_r(x)\subseteq U_i$ for some $i\in I$.
\end{lemma}
\textbf{Proof}: Before proving this, try to visualize the result!

We prove this lemma through contradiction. Assume that for some $r>0$ there exists an $x\in X$ (possibly depending on $r$) such that for each $i\in I$, $B_r(x)\not\subseteq U_i.$ Consider the sequence $\{x_n\}_n$ in $X$ such that $B_{1/n}(x_n) \not\subseteq U_i$ for all $i\in I.$

Given that $X$ is sequentially compact, $\{x_n\}$ has a convergent subsequence $\{x_{n_k}\}_k$. Let $x_{n_k}\to x\in X$. Given that $\{U_i\}$ is an open cover of $X$, there exists a $U_{i_0}$ such that $x\in U_{i_0}$. Given $U_{i_0}$ is open, it also follows that there exists an $r_0$ such that $B_{r_0}(x) \subseteq U_{i_0}$. Hence, choose $N$ large enough such that $d(x,x_N) < \frac{r_0}{2}$ and $\frac{1}{N}< \frac{r_0}{2}$. Then, if $y\in B_{1/N}(x_N)$, then 
\[d(x,y) \leq  d(x,x_N) + d(x_N, y) <r_0.\]
Therefore, $y\in B_{r_0}(x) \subseteq U_{i_0}$. Hence,
\[
B_{1/N}(x_N) \subseteq B_{r_0} (x) \subseteq U_{i_0}
\]
which is a contradiction. \qed 

We call this $r$ the \textbf{Lebesgue number} of the open cover of $X$, which is useful in other applications.

\begin{definition}
A metric space $X$ it \textbf{totally bounded} if, for every $\epsilon>0$, there exists $x_1,x_2,\dots, x_k\in X$ with $k$ finite such that $\{B_\epsilon(x_i) \mid 1 \leq i \leq k\}$ is an open cover of $X$.
\end{definition}

\begin{lemma}
A metric space $X$ being sequentially compact implies that $X$ is totally bounded.
\end{lemma}
\textbf{Proof}: Assume that $X$ is sequentially compact and not totally bounded. Therefore, there exists an $\epsilon>0$ such that $X$ cannot be covered by a collection of open sets of only finitely many $\epsilon$-balls. Hence, let $x_1 \in X$, $x_2 \in X\setminus B_\epsilon(x_1)$, then $x_3 \in X\setminus B_\epsilon(x_1) \setminus B_\epsilon(x_2)$ and so on. We know that there exists such $x_i$ by the previous statement. Hence, for all $i\neq j$, $d(x_i, x_j) \geq \epsilon$. Therefore, $\{x_n\}_n$ has no convergent subsequence as if there was a convergent subsequence it would be Cauchy, and the previous line shows that no subsequence of $\{x_n\}$ will be Cauchy. This is a contradiction to $X$ being sequentially compact. \qed 

\begin{theorem}
A metric space $X$ is (topologically) compact if and only if $X$ is sequentially compact.
\end{theorem}

\textbf{Proof}:  We first show that sequentially compact implies topologically compact, applying the two lemmas developed thus far.  Let $X$ be sequentially compact and let $\{U_i\}_{i\in I}$ be an open cover of $X$. By the Lebesgue number lemma, there exists an $r>0$ such that for each $x\in X$, $B_r(x) \subset U_i$ for some $i\in I$. Furthermore, by Lemma 5, $X$ is totally bounded. Hence, there exists $y_1,\dots, y_k\in X$ such that 
\[
X\subset B_r(y_1) \cup \dots \cup B_r(y_k).
\]
However, for each $i\in I$, we have $B_r(y_i) \subset U_{j(i)}$ for some $j(i) \in I$. (This notation just means for each $i$, there exists a $j\in I$ which depends on $i$ such that $B_r(y_i) \subseteq U_{j}$.) Thus, $\{U_{j(1)},\dots, U_{j(k)}\}$ is a finite subcover for $X$. Therefore, every open cover of $X$ has a finite subcover, and thus sequentially compact implies topologically compact. 

We now prove the other direction. Assume for the sake of contradiction there there exists a sequence $\{x_n\}_n$ in $X$ with no convergent subsequence. Notice that no term in the sequence can appear infinitely many times, as otherwise there would be a trivial subsequence of $\{x_n\}$. Hence, we assume without loss of generality that $x_i \neq x_j$ if $i\neq j$. Furthermore, notice then that for every $n$ there exists an $\epsilon_n>0$ such that $B_{\epsilon_n}(x_n)$ contains no other terms in the sequence. If this wasn't the case, then there would again be a convergent subsequence of $\{x_n\}_n$. Therefore, for each $i$, there exists an open ball $U_i$ centered at $x_i$ such that $x_j \notin U_i$ for all $i\neq j$.

Additionally, consider $U_0 = X\setminus \{x_n\mid n\in \NN\}$. $U_0$ is open, as $U_0^c = \{x_n \mid n\in \NN\}$ is closed (it contains all of it's limit points). Hence,
\[
U_0 \cup \{U_n\mid n\in \NN\}
\]
is an open cover of $X$. However, this open cover has no finite subcover as any finite collection of the cover will fail to include infinitely many terms from the sequence $\{x_n\}_n$. This is a contradiction, and thus topologically compact implies sequentially compact. \qed 

\begin{remark}
Notice that we technically could've used this proof in the previous lecture, but the Heine-Borel Theorem is so vastly important that I decided to do that proof before today's lecture.
\end{remark}

We will now start to look at some illuminating applications of compact sets to reach an even more powerful theorem.

\begin{recall}
Let $X,Y$ be metric spaces and $f:X\to Y$ be a continuous function. Then, for all $U$ open in $Y$, $f^{-1}(U)$ is open in $X$.
\end{recall}

\begin{theorem}
Let $X,Y$ be metric spaces and $f:X\to Y$ be continuous. Given $K\Subset X$, $f(K) \subset Y$ is compact.
\end{theorem}
\textbf{Proof}: Let $\{U_i\}_{i\in I}$ be an open cover of $f(K)$. Then, define $V_i = \{f^{-1}(U_i)\}_{i\in I}$, which is open as $f$ is continuous. Therefore, $\{f^{-1}(U_i)\}_{i\in I}$ is an open cover of $K$. Hence, there exists a finite subcover $\{V_{i_1}, \dots V_{i_k}\}$ of $K$ as $K$ is compact. Thus, $\{U_{i_1},\dots U_{i_k}\} = \{f(V_{i_1}), \dots, f(V_{i_k})\}$ is a finite subcover of $f(K)$. Therefore, $f(K)$ is compact. \qed 
\begin{corollary}
Let $X$ be a metric space and $K\Subset X$. Then, given a continuous function $f:X\to \R$, $f$ obtains a maximum and minimum finite value on $K$.
\end{corollary}
\textbf{Proof}: The proof follows from the previous theorem, a problem on PSET 2. \qed 

\begin{corollary}
Sometimes in particular we want to study bounded continuous functions, and the previous corollary gives us a nice property. Given a compact metric space $X$, every continuous function on $f$ is bounded.
\end{corollary}
\textbf{Proof}: Follows immediately. \qed 

\begin{theorem}[Cantor's Intersection Theorem]
If $K_1 \supset K_2\supset K_3\supset \dots$ is a decreasing sequence of nonempty sequentially compact subsets of $\R^n$, then $\bigcap_{i\geq 1}K_i$ is non-empty.
\end{theorem}
\textbf{Proof}: Choose a sequence $\{a_n\}_n$ such that $a_n \in K_n$ for each $n$. We know that there exists such an $a_n$ as each $K_n$ is nonempty. Then, $\{a_n\}_n$ is a sequence in $K_1$, and thus there exists a convergent subsequence $\{a_{n_k}\}_k$ such that $a_{n_k}\to a\in K_1$. Furthermore, $\{a_n\}_{n=2}^\infty$ is a sequence in $K_2$, and thus contains a a convergent subsequence. Therefore, $a\in K_2$. Continuing this process, we get that $a\in K_i$ for all $i$. Thus, $a\in \bigcap_{i\geq 1} K_i$. \qed 

\begin{definition}[Finite Intersection Property]
A collection of closed sets $\{C_i\}_i$ has the \textbf{finite intersection property} if every finite subcollection has a nonempty intersection.
\end{definition}

Given Lemma 3 and the Cantor Intersection Theorem, it is clear that there are some relations between compact sets, nonempty intersections of sets, and totally bounded sets. We hence show the following theorem.

\begin{theorem}
Given a metric space $(X,d)$, the following are equivalent. 
\begin{enumerate}[label = (\arabic*)]
    \item $X$ is compact.
    \item $X$ is sequentially compact.
    \item $X$ is Cauchy complete and totally bounded.
    \item Every collection of closed subsets of $X$ with the finite intersection property has a non-empty intersection.
\end{enumerate}
\end{theorem}
We have shown (1) $\iff$ (2), and thus we show  (1) $\iff$ (4) and (2) $\iff$ (3) to finish the proof.

\textbf{Proof}: (1) $\implies$ (4): Assume for the sake of contradiction that there exists a collection of closed subsets $\{C_i\}_{i\in I}$ with the finite intersection property such that $\bigcap_{i\in I} C_i = \emptyset$. Given $C_i$ is closed in $X$ for all $i$, $U_i = C_i^c$ is open in $X$ for each $i.$ Then,
\[
\bigcup_{i\in I} U_i = \bigcup_{i\in I} C_i^c = \left(\bigcap_{i\in I} C_i\right)^c = \emptyset^c = X.
\]
Hence, the $U_i$ cover $X$. Given $X$ is compact, there exists a finite subcover $\{U_{i_1}, \dots, U_{i_k}\}$ of $X.$ Thus, 
\[
    X = \bigcup_{n=1}^k U_{i_n} = \left(\bigcap_{n=1}^k U_{i_n}^c\right)^c = \left(\bigcap_{n=1}^k C_{i_n}\right)^c.
\]
Therefore, $\bigcap_{n=1}^k C_{i_n}= \emptyset$ which is a contradiction with the finite intersection property.

(4) $\implies$ (1): Suppose that $\{U_i\}_{i\in I}$ is an open cover of $X$, and let $C_i = U_i^c$ for each $i\in I$. Assume for the sake of contradiction that no finite subset of the $U_i$ covers $X$. We show that $C_i$ has the finite intersection property. Assume for the sake of contradiction that $\{C_{n_1}, \dots, C_{n_k}\}$ satisfies $C_{n_1} \cap \dots \cap C_{n_k} = \emptyset$. Then, 
\[
    \bigcup_{i=1}^k U_{n_i} = \left(\bigcap_{i=1}^k U_{n_i}^c\right)^c = \left(\bigcap_{i=1}^k C_{i_k}\right)^c = \emptyset^c = X.
\]
This is a contradiction with the assumption that no subset of the $U_i$ covers $X.$ Hence, $\{C_i\}_{i\in I}$ satisfies the finite intersection property. Therefore, $\{C_i\}_{i\in I}$ has non-empty intersection; i.e. $\bigcap_{i\in I} C_i \neq \emptyset$. Then, $\bigcup_{i\in I} U_i \neq X$, which is a contradiction to the $U_i$ being an open cover for $X$. Thus, every open cover of $X$ has a finite open subcover.

(2) $\implies$ (3): We have already shown that $X$ being sequentially compact implies totally bounded, and hence we only need show that sequentially compact implies Cauchy complete. Let $\{x_n\}$ be a Cauchy sequence in $X$. Given $\{x_n\}$ is a sequence in $X$, there exists a convergent subsequence $\{x_{n_k}\}$ in $X$ such that $x_{n_k} \to x \in X$. Let $\epsilon>0$, and choose $N$ such that $d(x_i, x_j) < \epsilon/2$ whenever $i,j\geq N$. Next, choose $n_k>N$ such that $d(x_{n_k}, x) <\epsilon/2$. Then, 
\[
d(x,x_N) \leq d(x,x_{n_k}) + d(x_{n_k},x_N) < \epsilon.
\]
Thus, $x_n \to x\in X$ as $n\to \infty$. Therefore, every Cauchy sequence in $X$ converges to a point in $X$. Hence, $X$ is Cauchy complete. 

(3) $\implies$ (2): This part of the proof is quite difficult. Consider a sequence $\{x_n\}_n$ in $X$. Given $X$ is totally bounded, for every $n\in \NN$, there exists a finite set of points $\{y_{1}^{(n)}, \dots, y_{r(n)}^{(n)}\}$ such that $X\subset B_{\frac 1n} (y_1^{(n)}) \cup \dots \cup B_{\frac 1n} (y_{r(n)}^{(n)})$. Define 
\[
S_n = \{y_1^{(n)}, \dots, y_{r(n)}^{(n)}\}.
\]
We want to find a convergent subsequence of $\{x_n\}_n$. We do so by construction. Given $S_1$ is finite, there exists a $y_{n(1)}^{(1)}\in S_1$ such that $B_1(y_{n(1)}^{(1)})$ contains infinitely many points from $\{x_n\}_n$. Choose $z_1$ from this ball. Then, given $S_2$ is finite, there is a $y_{n(2)}^{(2)}$ such that infinitely many points from $\{x_n\}_n$ are in $B_1(y_{n(1)}^{(1)}) \cap B_{1/2}(y_{n(2)}^{(2)})$. Choose $z_2$ from this set. Continue this procedure for each $k >1$, selecting a $z_k$ from $\bigcap_{i=1}^k B_{\frac 1k}(y_{n(k)}^{(k)})$. Then, we show $\{z_n\}_n$ is Cauchy. Let $\epsilon>0$. Then, there exists an $N\in \NN$ such that $\frac{1}{N} < \epsilon$. Hence, for all $n,m \geq N$, 
\[
d(z_n, z_m) <  \frac{1}{N}< \epsilon.
\]
Therefore, by the Cauchy completeness of $X$, $\{z_n\}$ converges to a point in $X$. \qed 

\begin{remark}
Where do we use the fact that each ball has infinitely many points? We do in fact use this property in the proof. Try to figure out how!
\end{remark}